\chapter{Preparación del host y controlador PCIe}
\label{chap:host-driver}

\section{Comprobaciones iniciales}

Antes de instalar el controlador, comprueba la versión del sistema, la arquitectura y el kernel.
En Linux puedes guardar esta información en la carpeta \code{evidence}:

\begin{lstlisting}[language=bash]
mkdir -p evidence
{
  cat /etc/os-release
  uname -r
  uname -m
  lspci -nn
} | tee evidence/host-inventory.txt
\end{lstlisting}

\code{uname -m} suele mostrar \code{x86\_64} en un PC y \code{aarch64} en una
Raspberry Pi 5 con un sistema de 64 bits. La versión del kernel será necesaria para
comprobar la compatibilidad del controlador.

\section{Habilitar PCIe en Raspberry Pi 5}

Esta sección solo es necesaria si utilizas una Raspberry Pi 5. Los dispositivos HAT+
compatibles pueden detectarse automáticamente, mientras que para un adaptador PCIe no HAT+
Raspberry Pi indica que debe habilitarse la interfaz PCIe en el archivo de configuración
\cite{rpi_pcie}.

En Raspberry Pi OS actual, abre:

\begin{lstlisting}[language=bash]
sudoedit /boot/firmware/config.txt
\end{lstlisting}

Si esa ruta no existe, comprueba \code{/boot/config.txt}. Añade:

\begin{lstlisting}[language=bash]
dtparam=pciex1
\end{lstlisting}

Guarda el archivo y reinicia:

\begin{lstlisting}[language=bash]
sudo reboot
\end{lstlisting}

Raspberry Pi 5 utiliza PCIe Gen2 por defecto. Aunque es posible habilitar Gen3,
Raspberry Pi advierte que la placa no está certificada para funcionar de forma estable
a esa velocidad \cite{rpi_pcie}. Para la puesta en marcha se mantiene Gen2.

\section{Comprobar la conexión PCIe}

Después del arranque, comprueba que el sistema detecta la tarjeta:

\begin{lstlisting}[language=bash]
lspci -nn
\end{lstlisting}

Debe aparecer un dispositivo PCIe correspondiente a la tarjeta. Si no aparece, revisa
el montaje, la alimentación y, en Raspberry Pi 5, que PCIe esté habilitado antes de
continuar con el controlador.

Si necesitas más información sobre el dispositivo detectado, utiliza:

\begin{lstlisting}[language=bash]
lspci -nnk
\end{lstlisting}

En este punto solo se está comprobando la conexión PCIe. Que la tarjeta aparezca en
\code{lspci} no significa todavía que el controlador Akida esté instalado.

\section{Instalar el controlador}

El controlador PCIe oficial de BrainChip se distribuye en el repositorio
\code{akida\_dw\_edma}. Antes de instalarlo, comprueba la versión del kernel:

\begin{lstlisting}[language=bash]
uname -r
\end{lstlisting}

La versión del controlador consultada para este manual admite kernels desde 5.4 hasta
6.8. En kernels 6.9 o posteriores el módulo no se compila \cite{brainchip_driver}.
Como esta compatibilidad puede cambiar, consulta también el README del
repositorio antes de instalarlo.

\paragraph{Kernel no compatible.}
Si la versión del kernel queda fuera del intervalo soportado por el controlador,
no continúes con la instalación hasta comprobar si BrainChip ha publicado una versión
compatible.


En Ubuntu, instala las dependencias indicadas por BrainChip:

\begin{lstlisting}[language=bash]
sudo apt update
sudo apt install dkms build-essential "linux-headers-$(uname -r)"
\end{lstlisting}

En otras distribuciones los nombres de los paquetes pueden cambiar, pero deben
instalarse DKMS, las herramientas de compilación y los headers correspondientes al
kernel en uso.

A continuación, clona el repositorio e instala el controlador:

\begin{lstlisting}[language=bash]
git clone https://github.com/Brainchip-Inc/akida_dw_edma.git
cd akida_dw_edma
git rev-parse HEAD
./install.sh
\end{lstlisting}

La salida de \code{git rev-parse HEAD} registra la revisión exacta del controlador
utilizada; también puedes usar una configuración propia.

\section{Comprobar el controlador}

Después de instalar el controlador, revisa el estado de DKMS y confirma que el dispositivo
Akida está disponible:

\begin{lstlisting}[language=bash]
dkms status
lsmod | grep -i akida
ls -l /dev/akida*
lspci -nnk
\end{lstlisting}

La instalación puede considerarse correcta cuando el dispositivo sigue apareciendo
en PCIe, el módulo está cargado y existe al menos un nodo \code{/dev/akida*}.

Si el dispositivo aparece en \code{lspci} pero falta \code{/dev/akida*}, consulta los
mensajes del kernel:

\begin{lstlisting}[language=bash]
sudo dmesg | grep -iE 'akida|edma|pcie' | tail -n 120
\end{lstlisting}

Si el nodo existe pero solo puede utilizarse como \code{root}, revisa sus permisos y
la regla udev instalada por el controlador antes de continuar.

Con el controlador funcionando, continúa con la configuración del entorno de software
del \cref{chap:software}.
